\section{Non-Interactive Zero-Knowledge Proofs}

Non-interactive zero-knowledge proof systems enable a prover to convince a
verifier that a statement is true without revealing any information beyond the
validity of the statement.

\paragraph{Interfaces.} A non-interactive zero-knowledge proof (NIZK) system
consists of the following algorithms:

\begin{itemize}

  \item $\crs \sample \NIZKgen(\secparam)$ : the setup algorithm outputs a
    common reference string $\crs$.

  \item $\proof \leftarrow \NIZKprove(\crs, \inst, \witness)$ : the proving
    algorithm takes as input a common reference string $\crs$, statement
    $\inst$, and witness $\witness$, and outputs a proof $\proof$.

  \item $ \{0, 1\} \leftarrow \NIZKverify(\crs, \inst, \proof)$ : the verification
    algorithm takes as input a common reference string $\crs$, statement
    $\inst$, and proof $\proof$, and outputs $0$ or $1$.

\end{itemize}

\paragraph{Correctness.} A non-interactive zero-knowledge proof system must
satisfy the following correctness requirement: for every common reference
string $\crs$ generated by $\NIZKgen(\secparam)$ and every statement-witness
pair $(\inst, \ \witness) \in \lang$, the following holds:
\begin{equation*}
\NIZKverify(\crs, \inst, \NIZKprove(\crs, \inst, \witness)) = 1.
\end{equation*}

\paragraph{Security.} The constructions presented in this thesis require
different security notions for non-interactive zero-knowledge proof systems,
including zero knowledge, simulation soundness, and simulation extractability.
Since these notions are construction dependent, we defer their formal
definitions to the chapters in which they are required.